%==============================================================================
%  PEER REVIEW RUBRIC -- Final Project Paper
%  One form for all reviewers. The SAME form is used whether you are reviewing a
%  paper in your own category or another category; the weighting is applied by
%  the teaching staff (cross-category reviews and same-category reviews count
%  differently toward the author's grade). Score each criterion from 1 to 5.
%  Compile on Overleaf with pdfLaTeX.
%==============================================================================
\documentclass[11pt]{article}

\usepackage[a4paper,margin=2cm]{geometry}
\usepackage{booktabs}
\usepackage{tabularx}
\usepackage{array}
\usepackage{enumitem}
\usepackage{xcolor}
\usepackage{titlesec}

\titlespacing*{\section}{0pt}{1.0ex}{0.6ex}
\setlength{\parindent}{0pt}
\renewcommand{\arraystretch}{1.35}

% score column: fixed narrow, centered
\newcolumntype{S}{>{\centering\arraybackslash}p{2.2cm}}

\begin{document}

\begin{center}
{\large\bfseries Final Project --- Peer Review Rubric}\\[2pt]
Parallel Programming, Spring 2026
\end{center}

\vspace{0.4em}

% --- header fields filled by the reviewer ---
\begin{tabularx}{\textwidth}{@{}lX lX@{}}
Reviewer name: & \hrulefill & Reviewer ID: & \hrulefill \\[4pt]
Paper / Author ID: & \hrulefill & Paper category (A / B): & \hrulefill \\
\end{tabularx}

\vspace{0.8em}

\textbf{How to score.} Rate each of the seven criteria from 1 to 5. You are
\emph{not} expected to re-verify the technical work yourself. Instead, judge
whether the authors have explained and justified it convincingly --- something
any classmate in this course can assess, regardless of sub-topic.

\vspace{0.4em}

\textbf{Scale.}\quad
\textbf{5} Excellent \,\textbar\;
\textbf{4} Good \,\textbar\;
\textbf{3} Adequate \,\textbar\;
\textbf{2} Weak \,\textbar\;
\textbf{1} Poor / missing

\vspace{0.8em}

\begin{tabularx}{\textwidth}{@{}p{0.4cm} >{\raggedright\arraybackslash}p{4.0cm} X S@{}}
\toprule
\textbf{\#} & \textbf{Criterion} & \textbf{What to assess} & \textbf{Score (1--5)} \\
\midrule

1 & Problem \& Motivation &
Is it clear what problem is addressed and why it matters --- even to a reader
outside this sub-topic? & \\

2 & Background Adequacy &
Is enough background and baseline given to understand the approach? & \\

3 & Methodology Clarity \& Justification &
Is the approach described clearly, with reasonable justification for the key
decisions? & \\

4 & Core Contribution &
Is the case for the core contribution convincing?\newline
\textit{Cat.\ A:} the optimization / design technique.\newline
\textit{Cat.\ B:} the prompt-engineering and agent-guidance strategy. & \\

5 & Experimental Design \& Reproducibility &
Are the environment, measurement method, and correctness check described well
enough that the results could be reproduced? & \\

6 & Results \& Depth of Analysis &
Do the results support the claims, and is the analysis deep?\newline
\textit{Cat.\ A:} discussion of limitations.\newline
\textit{Cat.\ B:} root-cause analysis of the failed cases. & \\

7 & Writing \& IEEE Format &
Structure, language, and adherence to the IEEE paper format. & \\

\midrule
\multicolumn{3}{r}{\textbf{Total (out of 35)}} & \\
\bottomrule
\end{tabularx}

\vspace{1.0em}

{\color{gray}\footnotesize
Note: All reviewers use this single form. Technical criteria are phrased as
``is it justified / explained convincingly'' rather than ``is it correct,'' so
they can be answered by any reviewer in the course. The teaching staff applies
the cross-category and same-category weighting when combining scores.}

\end{document}